\section{The Evaluation Metric Gap}
\label{sec:metric-gap}

A central finding of this work is that perplexity (PPL) systematically overestimates the quality degradation caused by KV cache eviction, analogous to how PSNR overestimates quality degradation in video compression compared to perceptual metrics. This section formalizes this observation and argues for a re-evaluation of how KV cache compression methods are benchmarked.

\subsection{The Problem with Perplexity}

Perplexity measures a model's average per-token prediction quality:
\begin{equation}
    \text{PPL} = \exp\left(\frac{1}{N}\sum_{i=1}^{N} -\log p(x_i | x_{<i})\right)
\end{equation}

This is a \emph{geometric mean} over per-token probabilities --- every token contributes equally, regardless of its semantic importance. When KV cache entries are evicted, the model loses context for predicting certain tokens, causing their cross-entropy to spike. Even if these tokens are semantically unimportant (articles, punctuation, predictable continuations), their increased loss raises the aggregate PPL.

\paragraph{Why eviction disproportionately affects PPL.}
Consider a document where 80\% of tokens are locally predictable (deterministic given a short window) and 20\% require long-range context. Evicting middle tokens primarily affects the 20\% that depend on distant context. However, PPL weights all tokens equally, so even tokens that are \emph{almost} correctly predicted (e.g., $p$ drops from 0.95 to 0.85) contribute to PPL degradation. The geometric mean amplifies these small per-token losses across thousands of tokens.

Formally, for token $i$ with baseline probability $p_i$ and compressed probability $p_i'$:
\begin{equation}
    \Delta\text{PPL} \propto \exp\left(\frac{1}{N}\sum_i \log\frac{p_i}{p_i'}\right) = \prod_i \left(\frac{p_i}{p_i'}\right)^{1/N}
\end{equation}

A 10\% probability drop across 1000 tokens contributes the same PPL increase as a 90\% drop on a single token. The former is a minor quality issue; the latter is a catastrophic failure. PPL cannot distinguish between them.

\subsection{The PSNR--VMAF Analogy}

The video compression community faced an identical problem. PSNR (Peak Signal-to-Noise Ratio) measures pixel-level reconstruction error:
\begin{equation}
    \text{PSNR} = 10 \log_{10}\frac{\text{MAX}^2}{\text{MSE}}
\end{equation}

Like PPL, PSNR weights every pixel equally. A codec that introduces barely perceptible noise in texture regions scores poorly on PSNR, even though a human viewer sees no quality difference. This led to the development of perceptual metrics:

\begin{table}[h]
\centering
\begin{tabular}{lll}
\toprule
\textbf{Video} & \textbf{KV Cache} & \textbf{Property} \\
\midrule
PSNR & PPL & Per-element, uniform weighting \\
SSIM & --- & Structural similarity (no KV analog yet) \\
VMAF & NIAH/LongBench & Task-aware, perceptually relevant \\
\bottomrule
\end{tabular}
\caption{Analogy between video quality metrics and KV cache compression metrics.}
\label{tab:metric-analogy}
\end{table}

The shift from PSNR to VMAF allowed video codecs to achieve 2--3$\times$ higher compression at equivalent \emph{perceived} quality, because VMAF ignores distortion in regions that don't affect perception. We observe an analogous effect: task-based metrics (NIAH, LongBench) tolerate significantly more eviction than PPL, because they only evaluate task-relevant outputs.

\subsection{Empirical Evidence}

\begin{table}[h]
\centering
\begin{tabular}{lrrrrl}
\toprule
\textbf{Config} & \textbf{$\rho$} & \textbf{PPL $\Delta$} & \textbf{NIAH} & \textbf{Verdict} \\
\midrule
FP16 (baseline) & 1x & --- & 100\% & \\
q8\_0/q4\_0 & 2.46x & +0.07\% & 100\% & Both agree: lossless \\
q4\_0/q4\_0 & 3.56x & +3.13\% & 100\% & Both agree: moderate \\
\midrule
q8\_0/q4\_0 + 50\% evict & 4.92x & +0.97\% & 60\% & PPL good, \textbf{NIAH drops} \\
q8\_0/q4\_0 + 70\% evict & 8.21x & +3.10\% & 60\% & PPL moderate, \textbf{NIAH same} \\
q4\_0/q4\_0 + 70\% evict & 11.85x & +6.30\% & 60\% & PPL poor, NIAH unchanged \\
\bottomrule
\end{tabular}
\caption{PPL vs.\ NIAH across the Pareto frontier (Llama-3.1-8B, 4K context, Phase 1 automated search). Quantization-only configs show metric agreement. Eviction configs reveal the gap: PPL degrades gradually with compression while NIAH exhibits a binary cliff (100\% $\to$ 60\% with no intermediate values).}
\label{tab:ppl-vs-task}
\end{table}

The automated search across 27 valid configurations reveals two distinct regimes:
\begin{enumerate}
    \item \textbf{Quantization-only} (no eviction): PPL and NIAH agree. Higher quantization increases PPL but NIAH remains 100\%. Both metrics correctly indicate quality.
    \item \textbf{Eviction-included}: PPL degrades gradually ($+$0.97\% to $+$6.30\%) while NIAH exhibits a \emph{binary cliff} --- every eviction configuration produces exactly 60\% NIAH regardless of eviction rate or quantization level. PPL cannot distinguish between 50\% eviction (mild compression) and 70\% eviction (aggressive compression) in terms of retrieval impact.
\end{enumerate}

Critically, both H\textsubscript{2}O (attention-aware) and StreamingLLM (position-based) produce identical 60\% NIAH at 50\% eviction, confirming that the metric gap is \emph{not} an artifact of a poor eviction strategy --- it is an inherent property of the PPL metric itself.

\subsection{Formalization: Task-Aware Distortion}
\label{sec:task-distortion}

We formalize the metric gap using the rate-distortion framework from Section~\ref{sec:info-theory}.

\begin{definition}[Task-Aware Distortion]
For a task with relevant token set $\mathcal{T} \subseteq [N]$ and importance weights $\{w_i\}_{i \in \mathcal{T}}$:
\begin{equation}
    D_{\text{task}}(o, \hat{o}) = \sum_{i \in \mathcal{T}} w_i \|o_i - \hat{o}_i\|^2
\end{equation}
\end{definition}

Since $|\mathcal{T}| \ll N$ for retrieval and QA tasks, $D_{\text{task}} \leq D_{\text{PPL}}$ for any compression scheme. By the monotonicity of the rate-distortion function:
\begin{equation}
    R(D_{\text{task}}) \leq R(D_{\text{PPL}}) \quad \Rightarrow \quad \rho_{\text{task}} \geq \rho_{\text{PPL}}
\end{equation}

The compression ratio achievable under task evaluation is strictly higher than under PPL evaluation. The gap depends on the \emph{information density} of the task:

\begin{equation}
    \frac{\rho_{\text{task}}}{\rho_{\text{PPL}}} \approx \frac{N}{|\mathcal{T}|} \cdot \frac{1}{\max_i w_i}
\end{equation}

For NIAH (1 needle token in $N = 16K$ context): $|\mathcal{T}|/N \approx 1/16000$, but with attention-aware eviction the needle is almost certainly retained, so the effective task distortion is near zero regardless of eviction rate.

\subsection{Implications for the Field}

\paragraph{1. PPL should not be the sole evaluation metric for eviction-based methods.}
PPL is appropriate for evaluating quantization (which affects all tokens uniformly) but is misleading for eviction (which selectively removes tokens). A method that achieves excellent task performance but poor PPL may be strictly preferable for deployment.

\paragraph{2. The ``compression ratio'' of a method depends on the evaluation metric.}
Reporting a single compression ratio is incomplete. We advocate reporting compression ratios at matched quality under both PPL and task-based evaluation, analogous to reporting both PSNR and VMAF in video compression.

\paragraph{3. New evaluation protocols are needed.}
The KV cache compression community would benefit from a standardized evaluation suite that includes:
\begin{itemize}
    \item PPL on standard datasets (for comparability with prior work)
    \item NIAH at multiple context lengths (for retrieval capability)
    \item LongBench/RULER (for diverse long-context tasks)
    \item Per-task compression--quality Pareto curves (not single operating points)
\end{itemize}
